\begingroup
\renewcommand{\arraystretch}{1.45}
\fancyhead[L]{\small ARITHMETISCH-GEOMETRISCHES MITTEL.}
\Needspace{28\baselineskip}
\AGMArticle{23}
[Bildet man denselben Algorithmus wie zuvor für $\alpha,\allowbreak \beta,\allowbreak \gamma,\allowbreak \delta$ nun für $A,\allowbreak B,\allowbreak C,\allowbreak D$ und setzt $A=\gamma,\allowbreak \ B=\delta$, so gilt offenbar für jedes $n$: $A^{\mathrm n}=\gamma^{\mathrm n},\allowbreak \ B^{\mathrm n}=\delta^{\mathrm n},\allowbreak \ C^{\mathrm n}=\alpha^{\mathrm n},\allowbreak \ D^{\mathrm n}=\beta^{\mathrm n}$ und somit für die Grenzwerte $K,\allowbreak H$, die
\SourceLeaf{405}
$x,\allowbreak \eta$ entsprechen: $\frac{x}k=\frac Kk\cdot y^{\frac12}\cdot H,\allowbreak \ \frac{x}k y^{\frac12}\eta=\frac Kk$; demnach gelten die Funktionalgleichungen
\Needspace{14\baselineskip}
\[
\begin{aligned}\dfrac{R(y,\dfrac1{y\eta})}P
&=\dfrac{S(y,\dfrac1{y\eta})}Q
\\[7pt]&=\dfrac{P(y,\dfrac1{y\eta})}R
\\[7pt]&=\dfrac{Q(y,\dfrac1{y\eta})}S
\\[7pt]&=y^{-\frac14}\eta^{-\frac12}\end{aligned}
\]

Setzt man dagegen $A=\beta,\allowbreak \ B=\alpha$, so wird $C=-\delta,\allowbreak \ D=-\gamma$, und für jedes $n$, das gleich oder größer als eins ist: $A^{\mathrm n}=\alpha^{\mathrm n},\allowbreak \ B^{\mathrm n}=\beta^{\mathrm n},\allowbreak \ C^{\mathrm n}=\gamma^{\mathrm n},\allowbreak \ D=\delta^{\mathrm n}$; man erhält also:
\Needspace{14\baselineskip}
\[
\begin{aligned}\dfrac{Q(y,-\eta)}P
&=\dfrac{P(y,-\eta)}Q
\\[7pt]&=-\dfrac{iS(y,-\eta)}R
\\[7pt]&=-\dfrac{iR(y,-\eta)}S\\[7pt]&=1\end{aligned}
\]

\Needspace{22\baselineskip}
Setzt man schließlich $2\surd {A''}=\surd \alpha+\surd \beta,\allowbreak \ 2\surd {C''}=\surd \alpha-\surd \beta$, so gilt für jedes $n\geqq1$
\[
\begin{gathered}
2A^{\mathrm n+1}=\surd {a^{\mathrm n}\alpha^{\mathrm n}}+\surd {b^{\mathrm n}\beta^{\mathrm n}},\\[7pt]
2C^{\mathrm n+1}=\surd {a^{\mathrm n}\alpha^{\mathrm n}}-\surd {b^{\mathrm n}\beta^{\mathrm n}}\\[7pt]
2B^{\mathrm n+1}=\surd {b^{\mathrm n}\alpha^{\mathrm n}}+\surd {a^{\mathrm n}\beta^{\mathrm n}},\\[7pt]
2D^{\mathrm n+1}=\surd {b^{\mathrm n}\alpha^{\mathrm n}}-\surd {a^{\mathrm n}\beta^{\mathrm n}}\\[7pt]
2\surd {A^{\mathrm n+2}}=\surd {\alpha^{\mathrm n}}+\surd {\beta^{\mathrm n}},\\[7pt]
2\surd {C^{\mathrm n+2}}=\surd {\alpha^{\mathrm n}}-\surd {\beta^{\mathrm n}},\\[7pt]
4A^{\mathrm n+1}C^{\mathrm n+1}=c^{\mathrm n}\gamma^{\mathrm n}\\[7pt]
\left(\dfrac Kk\right)^4=\dfrac{x}k,\\[7pt] HH=\eta
\end{gathered}
\]
also:
\Needspace{12\baselineskip}
\[
\begin{aligned}
2P(y^4,\eta\eta)&=P+Q,\\[6pt]2R(y^4,\eta\eta)&=P-Q\\[6pt]
2P(yy,\eta)^2&=pP+qQ,\\[6pt]2R(yy,\eta)^2&=pP-qQ\\[6pt]
2Q(yy,\eta)^2&=qP+pQ.\\[6pt]2S(yy,\eta)^2&=qP-pQ
\end{aligned}
\]
\Needspace{12\baselineskip}
\[
\begin{aligned}
p(\surd  y)\cdot P(\surd  y,\eta)&=PP+RR,\\[6pt]
q(\surd  y)\cdot P(\surd  y,\eta)&=QQ-SS,\\[6pt]
r(\surd  y)\cdot R(\surd  y,\eta&=2PR\\[6pt]
q(\surd  y)\cdot Q(\surd  y,\eta)&=PP-RR,\\[6pt]
p(\surd  y)\cdot Q(\surd  y,\eta)&=QQ+SS,\\[6pt]
r(\surd  y)\cdot S(\surd  y,\eta&=2QS
\end{aligned}
\]
wobei sich wiederum diejenigen Funktionen, denen keine Argumente beigefügt sind, auf die einfachen $y$ und $\eta$ beziehen.

Der so erhaltene neue Algorithmus, bei dem man von den Funktionen $P,\allowbreak Q,\allowbreak R,\allowbreak S$ mit den Argumenten $y,\allowbreak \eta$ zu den Funktionen mit den Argumenten $yy,\allowbreak \eta$ übergeht, ist offenbar der von \textsc{Gauss} in Artikel 16 der \textit{Determinatio attractionis quam in punctum quodvis positionis datae exerceret planeta etc.}\ verwendete.

Die Beziehungen zwischen den Funktionen, die sich auf die beiden beliebigen Werte $\xi,\allowbreak \eta$ des zweiten Arguments beziehen, und denjenigen Funktionen mit den zweiten Argumenten $\xi\eta$ und $\frac\xi\eta$ lassen sich auf verschiedene Weise mit
\SourceLeaf{406}
Hilfe des für $\alpha,\allowbreak \beta$ aufgestellten Algorithmus herleiten. Die Methode mit den am wenigsten umfangreichen Entwicklungen scheint diejenige zu sein, die sich auf Funktionen bezieht, bei denen neben den zweiten Argumenten $\xi\eta$ und $\frac{\xi}{\eta}$ als erstes Argument das Quadrat des ersten Arguments der Funktionen auftritt, die $\xi,\allowbreak \eta$ als zweites Argument haben.]

\Needspace{22\baselineskip}
\AGMArticle{24}
[Beim Algorithmus des arithmetisch-geometrischen Mittels ergeben die durch Rückwärtsverlängerung entstehenden Glieder mithilfe der Gleichungen
\Needspace{12\baselineskip}
\[
\begin{aligned}\dfrac{a_{\mathrm n}}{{}^{\mathrm n}\!a}&=\dfrac{c_{\mathrm n}}{{}^{\mathrm n}\!c}\\[7pt]&=\dfrac{b_{\mathrm n}}{{}^{\mathrm n}\!b}\\[7pt]&=\dfrac{1}{2^{\mathrm n}}\end{aligned}
\]
\Needspace{28\baselineskip}
die Glieder $a_{\mathrm n},\allowbreak c_{\mathrm n},\allowbreak b_{\mathrm n}$, die ebenso von $a,\allowbreak c$ abhängen wie $a^{\mathrm n},\allowbreak b^{\mathrm n},\allowbreak c^{\mathrm n}$ von $a,\allowbreak b$. Vertauscht man dementsprechend beim kombinierten Algorithmus gleichzeitig $b$ mit $c$ und $\beta$ mit $\gamma$ und lässt $\alpha$ unverändert, so wechselt $\delta$ nur sein Vorzeichen; wir erhalten also einen Algorithmus, der ebenso von $a,\allowbreak c,\allowbreak \alpha,\allowbreak \gamma$ abhängt wie der bisher betrachtete von $a,\allowbreak b,\allowbreak \alpha,\allowbreak \beta$, wenn wir setzen
\[
\begin{gathered}
\alpha_0=\alpha,\\[7pt]\gamma_0=\beta,\\[7pt]\beta_0=\gamma,\\[7pt]\delta_0=-\delta\\[8pt]
\begin{aligned}
\alpha_1&=\dfrac{1}{a_1}\left(\dfrac{\alpha_0+\gamma_0}{2}\right)^2
\\[6pt]&=\dfrac{1}{b_1}\left(\dfrac{\beta_0-\delta_0}{2}\right)^2,
\\[7pt]\gamma_1&=\dfrac{1}{c_1}\alpha_0\gamma_0\\[8pt]
\beta_1&=\dfrac{1}{b_1}\left(\dfrac{\alpha_0-\gamma_0}{2}\right)^2
\\[6pt]&=\dfrac{1}{a_1}\left(\dfrac{\beta_0+\delta_0}{2}\right)^2,
\\[7pt]\delta_1&=\dfrac{1}{c_1}\beta_0\delta_0
\end{aligned}
\end{gathered}
\]
\Needspace{20\baselineskip}
und so weiter bis zu den Grenzwerten:
\[
\begin{gathered}
\lim a_{\mathrm n}=\lim c_{\mathrm n}=\mathrm{M}(a,c)=l\\[7pt]
\lim{}^{2^{\mathrm n}}\!\surd {\dfrac{b_{\mathrm n}}{4l}}=e^{-\dfrac{\pi}{2}\dfrac{\mathrm{M}(a,c)}{\mathrm{M}(a,b)}}=\surd  z\\[7pt]
\lim{}^{2^{\mathrm n}}\!\surd {\dfrac{\alpha_{\mathrm n}}{l}}=\lim{}^{2^{\mathrm n}}\!\surd {\dfrac{\gamma_{\mathrm n}}{l}}=\dfrac{\lambda}{l}\\[7pt]
{}^{2^{n-1}}\!\surd {\left[\surd {\dfrac{\beta_{\mathrm n}}{4l}}\mathbin{\pm}\surd {\dfrac{\delta_{\mathrm n}}{4l}}\right]}=\dfrac{\lambda}{l}z^{\frac12}\zeta^{\pm1}
\end{gathered}
\]
Es sind also $\alpha,\allowbreak \gamma,\allowbreak \beta,\allowbreak -\delta$ ebensolche Funktionen von $l,\allowbreak z,\allowbreak \lambda,\allowbreak \zeta$ wie $\alpha,\allowbreak \beta,\allowbreak \gamma,\allowbreak \delta$ von $k,\allowbreak y,\allowbreak x,\allowbreak \eta$. Da $\frac{l}{\lambda}$ allein von $y$ und $\eta$ abhängt, und zwar auf dieselbe Weise wie $\frac{\lambda}{l}$ von $\eta$ und $y$, muss für $\frac{\log y}{\pi}=\frac{\pi}{\log z}$,

\SourceLeaf{407}

\Needspace{22\baselineskip}
\[
\begin{aligned}\dfrac{P(z,\zeta)}{P(y,\eta)}&=\dfrac{R(z,\zeta)}{Q(y,\eta)}\\[7pt]&=\dfrac{Q(z,\zeta)}{R(y,\eta)}\\[7pt]&=\dfrac{iS(z,\zeta)}{S(y,\eta)}\\[7pt]&=T(y,\eta)\\[7pt]&=\dfrac{1}{T(z,\zeta)}\end{aligned}
\]
\Needspace{10\baselineskip}
gelten. Dabei ist außer $T$ auch noch $\zeta$ als Funktion von $y,\allowbreak \eta$ zu bestimmen, und zwar $\zeta$ so, dass $y,\allowbreak \eta,\allowbreak \zeta$ der Reihe nach mit $z,\allowbreak \zeta,\allowbreak \eta$ vertauscht werden können, also so, dass
\[
\begin{aligned}\dfrac{\log y}{\pi}&=\dfrac{\pi}{\log z}\\[7pt]&=\varphi[(\log\eta)^2,(\log\zeta)^2]\\[7pt]&=\dfrac{1}{\varphi[(\log\zeta)^2,(\log\eta)^2]}\end{aligned}
\]
gilt.

\Needspace{24\baselineskip}
Für $P,\allowbreak Q,\allowbreak R$ als Quadratwurzeln aus $\frac{\alpha}{x},\allowbreak \frac{\beta}{x},\allowbreak \frac{\gamma}{x}$ sind hier gleiche Vorzeichen gewählt, weil nach Artikel 18 $\eta$ und $\zeta$ zugleich den Wert $1$ annehmen und somit nach den Artikeln 22 und 17
\[
\begin{aligned}\dfrac{P(z,1)}{P(y,1)}&=\dfrac{R(z,1)}{Q(y,1)}\\[7pt]&=\dfrac{Q(z,1)}{R(y,1)}\\[7pt]&=T(y,1)\\[7pt]&=\dfrac{1}{T(z,1)}\\[7pt]&=\dfrac{pz}{py}\\[7pt]&=\dfrac{rz}{qy}\\[7pt]&=\dfrac{qz}{ry}\\[7pt]&=\surd {\dfrac{-\log y}{\pi}}\\[7pt]&=\surd {\dfrac{\pi}{-\log z}}\end{aligned}
\]
gilt. Dabei ist dasjenige Vorzeichen der Quadratwurzel zu nehmen, für das der reelle Teil positiv ist.

Verwandelt man die zweiten Argumente $\zeta,\allowbreak \eta$ in ihre Kehrwerte, so bleiben $P,\allowbreak Q,\allowbreak R$ unverändert und $S$ wechselt nur sein Vorzeichen; es gilt also
\Needspace{5\baselineskip}
\[
T(y,\eta)=T(y,\tfrac{1}{\eta})
\]
und eine entsprechende Bedingung gilt für $y$ als Funktion von $\eta,\allowbreak \zeta$, was durch die oben angegebene Form der $\varphi$-Funktion angedeutet sein soll.

\Needspace{18\baselineskip}
Transformiert man in der Gleichung, durch die $T$ eingeführt wurde, die $P,\allowbreak Q,\allowbreak R,\allowbreak S$ so, dass die Argumente $z,\allowbreak \zeta,\allowbreak y,\allowbreak \eta$ einmal in $z,\allowbreak -\zeta,\allowbreak y,\allowbreak \frac{1}{y\eta}$ und dann in $\surd  z,\allowbreak \zeta,\allowbreak yy,\allowbreak \eta\eta$ übergehen, und berücksichtigt man, dass
\[
\begin{aligned}\dfrac{p\surd  z}{pyy}&=\dfrac{r\surd  z}{qyy}\\[7pt]&=\dfrac{q\surd  z}{ryy}\\[7pt]&=\surd {\dfrac{-\log yy}{\pi}}\\[7pt]&=\surd {\dfrac{\pi}{-\log\surd  z}}\end{aligned}
\]
gilt, so erhält man
\Needspace{24\baselineskip}
\[
\begin{gathered}\begin{aligned}\dfrac{Q(z,-\zeta)}{R(y,\dfrac{1}{y\eta})}&=\dfrac{iS(z,-\zeta)}{S(y,\dfrac{1}{y\eta})}\\[7pt]&=\dfrac{P(z,-\zeta)}{R(y,\dfrac{1}{y\eta})}\\[7pt]&=\dfrac{R(z,-\zeta)}{Q(y,\dfrac{1}{y\eta})}\\[7pt]&=y^{\frac14}\eta^{\frac12}T(y,\eta)\\[7pt]&=T(y,\tfrac{1}{y\eta})\end{aligned}\\[10pt]\begin{aligned}\dfrac{P(\surd  z,\zeta)}{P(yy,\eta\eta)}&=\dfrac{R(\surd  z,\zeta)}{Q(yy,\eta\eta)}\\[7pt]&=\dfrac{Q(\surd  z,\zeta)}{R(yy,\eta\eta)}\\[7pt]&=\dfrac{iS(\surd  z,\zeta)}{S(yy,\eta\eta)}\\[7pt]&=2\surd {\dfrac{\pi}{-\log yy}}\cdot T(y,\eta)^2\\[7pt]&=T(yy,\eta\eta)\end{aligned}\end{gathered}
\]

\SourceLeaf{408}

\Needspace{22\baselineskip}
\[
\begin{gathered}\begin{aligned}\dfrac{\log y}{\pi}&=\dfrac{\pi}{\log z}\\[7pt]&=\varphi[(\log\eta)^2,(\log\zeta)^2]\\[7pt]&=\varphi[(\log\eta)^2,(\pi i+\log\zeta)^2]\end{aligned}\\[10pt]\begin{aligned}\dfrac{\log yy}{\pi}&=\dfrac{\pi}{\log\surd  z}\\[7pt]&=2\varphi[(\log\eta)^2,\log\zeta)^2]\\[7pt]&=\varphi[(\log\eta\eta)^2,(\log\zeta)^2]\end{aligned}\end{gathered}
\]
und daher
\Needspace{8\baselineskip}
\[
\begin{gathered}
T(y,\eta)=\surd {\dfrac{-\log y}{\pi}}\cdot e^{\frac{(\log\eta)^2}{4\log y}}\\[7pt]
\dfrac{\log y}{\pi}=\dfrac{\pi}{\log z}=\mathbin{\pm}i\cdot\dfrac{\log\eta}{\log\zeta}
\end{gathered}
\]
Dabei ist das Vorzeichen der Quadratwurzel so zu wählen, dass ihr reeller Teil positiv ist, das Vorzeichen von $\mathbin{\pm}i\frac{\log\eta}{\log\zeta}$ hingegen so, dass der reelle Teil dieses Ausdrucks negativ ist.]

\clearpage
\endgroup
